\begin{center}\begin{tikzpicture}[font=\small,>=stealth,line join=round]
\draw[gray!40,rounded corners=3pt] (0,0) rectangle (14.2,8.3);\node at (7.1,7.700000000000001) {例1：经过(4,2)，确定为平方根函数};
\begin{scope}[shift={(1.676923076923077,2.243548387096774)},x=1.1538461538461537cm,y=1.4838709677419353cm]\draw[->,gray] (-0.5,0)--(4.7,0) node[right]{$x$};\draw[->,gray] (0,-0.4)--(0,2.7) node[above]{$y$};\draw[violet,thick,solid,domain=0:4.2,samples=180] plot(\x,{sqrt(\x)});\fill[orange!85!black] (0,0) circle(2.2pt);\fill[orange!85!black] (4,2) circle(2.2pt);\fill[orange!85!black] (0.5,0.7071067811865476) circle(2.2pt);\fill[orange!85!black] (0.8,0.8944271909999159) circle(2.2pt);\node[above left=5pt] at (4,2){$(4,2)$};\node[below left] at (0,0){$O$};\end{scope}\node[anchor=west,align=left,text width=5.0cm] at (8.4,4.0) {$4^\alpha=2\Rightarrow\alpha=\frac12$\\[9pt]$f(x)=\sqrt{x}$\\[9pt]点$(4,2)$确定指数\\[9pt]$0.5<0.8$\\[9pt]所以$f(0.5)<f(0.8)$};
\node[text width=12.6cm,align=center] at (7.1,.65) {曲线完整定义域为非负数；两个近处橙点按横坐标顺序排列，函数值也同序。};
\end{tikzpicture}\end{center}
